\section{An explicit Boolean-to-algebraic edge with its encoding}
\label{sec:boolean}
The inclusion $\mathrm P\subseteq\mathrm{NP}$ follows directly:
a deterministic machine is a nondeterministic machine with one
available move at each step and the same running time.
The reverse inclusion is the equality question.
We retain a concrete edge by arithmetizing a three-literal conjunctive
formula's verifier.

For $n,m\ge1$, use $6mn$ commuting variables $E_{j,k,i,\sigma}$,
where $1\le j\le m$, $1\le k\le3$, $1\le i\le n$,
$\sigma\in\{+,-\}$, and auxiliary variables $b_1,\ldots,b_n$.
Over $\mathbb Z$ define
\begin{align*}
 L_{j,k}(E,b)&=\sum_{i=1}^n
   (E_{j,k,i,+}b_i+E_{j,k,i,-}(1-b_i)),\\
 V_{n,m}(E,b)&=\prod_{j=1}^m
   \left(1-\prod_{k=1}^3(1-L_{j,k}(E,b))\right),\\
 Q_{n,m}(E)&=\sum_{b\in\{0,1\}^n}V_{n,m}(E,b).
\end{align*}
A formula $\varphi$ with $n$ explicitly named variables and $m$
ordered three-literal clauses gives the zero--one array $E(\varphi)$
with one $1$ in each slot $(j,k)$, at its literal's index and sign.
Repeated literals are allowed. Input consists of explicitly named
variables and ordered signed binary indices with delimiters.
Writing $E(\varphi)$ uses $O(mn)$ bits and operations.
This convention does not allow exponentially many unused variables
to be specified solely by a binary upper bound: alternatively relabel
the $r$ occurring variables and record the unused count separately.
The projection from all $n$ Boolean variables to those $r$ coordinates
has exactly $2^{n-r}$ elements in each fibre and preserves the verifier.
Thus the original count is $2^{n-r}$ times the relabelled count.
A binary specification of an enormous unused count would require this
factored output format; writing all bits of the expanded count need
not be polynomial in that shorter input encoding.

\begin{proposition}[Exact counting reduction]\label{prop:counting}
The integer $Q_{n,m}(E(\varphi))$ is exactly the number of satisfying
assignments of $\varphi$. Deciding satisfiability of these formulas
therefore reduces by this polynomial array map and a nonzero test
to evaluation of this family on the specified arrays.
\end{proposition}
\begin{proof}
For Boolean $b$, each selected $L_{j,k}$ is exactly the literal's
truth value, $b_i$ or $1-b_i$.
The product of their three complements is $1$ exactly when the
clause is false. Its complement is its truth value.
The outer product is $1$ exactly when all clauses hold, and $0$
otherwise. Summation proves the integer equality and the bound
$0\le Q_{n,m}(E(\varphi))\le2^n$.
The stated array construction and zero test give the reduction.
\end{proof}

The displayed construction gives $V_{n,m}$ a uniform division-free
arithmetic circuit of $O(mn)$ gates, constants $0,1$, total degree
at most $6m$ in $(E,b)$ and degree at most $3m$ in $E$.
Each literal has degree at most two, each clause at most six, and
the products give these bounds. The loops over the displayed indices
construct the circuit in polynomial time. To obtain the usual one-index
family, assign
$t=(n+m-2)(n+m-1)/2+n$. The consecutive diagonals of pairs $(n,m)$
show this is a bijection from positive pairs to positive integers.
Trying consecutive diagonals through $n+m\le t+1$ decodes it in
polynomial time, and $n,m\le t$, so all variable, degree and gate bounds
remain polynomial in $t$. If a definition requires exactly $t$ witness
bits, multiply the verifier by $\prod_{i=n+1}^{t}(1-b_i)$ before
summing; only the all-zero added bits contribute, leaving $Q_{n,m}$
unchanged. By the definition of $\mathrm{VNP}$
as Boolean exponential summation of a polynomial-size,
polynomially bounded-degree arithmetic circuit family, this proves
$Q_{n,m}\in\mathrm{VNP}_{\mathbb Q}$ by construction.

There is also an explicit bit interface for a supplied, topologically
ordered, fan-in-two division-free circuit $C$ over $+,-,\times$, with
signed binary integer constants. On input $E(\varphi)$, evaluate its gates
modulo $M=2^{n+1}$. For $L$ gates and total constant-bit length $B_C$,
this takes $O(B_C+L(n+1)^2)$ arithmetic bit operations, plus graph and
input reading and the chosen machine's memory-access overhead,
by binary addition and school multiplication. Reduction to low
$n+1$ bits handles products, long constants and signs.
Induction over gates proves exact agreement with the integer value
modulo $M$. If the supplied circuit computes $Q_{n,m}$, its residue
is its exact value on this input since $0\le Q\le2^n<M$.
This gives the precise evaluation reduction without a bound on
intermediate integer magnitudes.

No polynomial-size circuits for $Q_{n,m}$ have been constructed here.
Polynomial-time generation of their full binary descriptions would
make the displayed evaluation uniform and bound both gate count and
constant-bit length. Descriptions supplied as size-dependent advice
give a polynomial-advice interface only when their total binary
description length is polynomial.
Arbitrary complex constants do not admit this binary modular
operation without additional input data.
The retained inverse-fibre selector polynomial already has the
explicit small determinant proved earlier; the present counting
polynomial has the different, fully specified joint verifier product
and Boolean summation. The next reduction must transport those
joint constraints with their parameter growth. None is erased by
calling both objects polynomial families.
